\begin{obeactivity}
\textbf{Aktivitas 8.1: Enkripsi dan Dekripsi Teks dengan Python}\\
1. Pasang pustaka PyCryptodome dengan perintah \texttt{pip install pycryptodome}.
2. Jalankan program pada Listing~\ref{lst:aes-teks}.
3. Masukkan pesan \texttt{Besok tunggu saya di Stasiun Bandung pukul 17.30} dan
   kunci \texttt{abcdefghij123456}.
4. Catat cipherteks yang muncul. Jalankan program sekali lagi dengan pesan dan
   kunci yang sama, lalu bandingkan cipherteksnya.
5. Perhatikan bahwa cipherteks dapat berbeda setiap penjalanan karena
   \textit{initialization vector} (IV) acak. Pastikan bahwa hasil dekripsi
   selalu sama dengan pesan semula.
\end{obeactivity}

\begin{lstlisting}[
    language=Python,
    float=htbp,
    basicstyle=\ttfamily\scriptsize,
    caption={Program enkripsi dan dekripsi teks dengan AES-128 mode CBC},
    label={lst:aes-teks}
]
from Crypto.Cipher import AES
from Crypto.Random import get_random_bytes
from Crypto.Util.Padding import pad, unpad
import base64
BLOCK_SIZE = 16

def enkripsi_aes(plainteks, kunci):
    if isinstance(plainteks, str):
        plainteks = plainteks.encode()
    kunci = kunci.encode()
    iv = get_random_bytes(BLOCK_SIZE)
    cipher = AES.new(kunci, AES.MODE_CBC, iv)
    data = pad(plainteks, BLOCK_SIZE)
    hasil = cipher.encrypt(data)
    return base64.b64encode(iv + hasil)

def dekripsi_aes(cipherteks, kunci):
    kunci = kunci.encode()
    data = base64.b64decode(cipherteks)
    iv = data[:BLOCK_SIZE]
    cipher = AES.new(kunci, AES.MODE_CBC, iv)
    hasil = cipher.decrypt(data[BLOCK_SIZE:])
    return unpad(hasil, BLOCK_SIZE)
pesan = input("Pesan: ")
kunci = input("Kunci (16 karakter): ")
cipherteks = enkripsi_aes(pesan, kunci)
print("Cipherteks:", cipherteks)
plainteks = dekripsi_aes(cipherteks, kunci)
print("Plainteks:", plainteks.decode())
\end{lstlisting}
% Koreksi: slide halaman 47 mencetak plainteks[16:]. Fungsi dekripsi sudah
% melakukan unpad, sehingga potongan 16 byte itu membuang awal pesan. Baris
% di atas memakai plainteks.decode() agar seluruh pesan tercetak.

\begin{obeactivity}
\textbf{Aktivitas 8.2: Enkripsi dan Dekripsi Berkas Biner}\\
1. Siapkan sebuah berkas, misalnya berkas PDF. Catat ukuran dan nilai
   \textit{checksum}-nya.
2. Tambahkan fungsi pada Listing~\ref{lst:aes-berkas} ke program Aktivitas 8.1.
3. Enkripsikan berkas PDF itu menjadi berkas \texttt{.enc}, lalu dekripsikan
   kembali menjadi berkas PDF baru.
4. Bandingkan nilai \textit{checksum} berkas asli dan berkas hasil dekripsi.
   Keduanya harus sama.
5. Jelaskan mengapa ukuran berkas hasil enkripsi lebih besar daripada berkas
   aslinya.
\end{obeactivity}

\begin{lstlisting}[
    language=Python,
    float=htbp,
    basicstyle=\ttfamily\scriptsize,
    caption={Fungsi enkripsi dan dekripsi berkas biner dengan AES-128},
    label={lst:aes-berkas}
]
def enkripsi_berkas(nama_masukan, nama_keluaran, kunci):
    with open(nama_masukan, "rb") as f:
        data = f.read()
    hasil = enkripsi_aes(data, kunci)
    with open(nama_keluaran, "wb") as f:
        f.write(hasil)

def dekripsi_berkas(nama_masukan, nama_keluaran, kunci):
    with open(nama_masukan, "rb") as f:
        data = f.read()
    hasil = dekripsi_aes(data, kunci)
    with open(nama_keluaran, "wb") as f:
        f.write(hasil)

kunci = "abcdefghij123456"
enkripsi_berkas("dokumen.pdf", "dokumen.enc", kunci)
dekripsi_berkas("dokumen.enc", "dokumen-hasil.pdf", kunci)
\end{lstlisting}
